\chapter{Results and Verification}

\section{Result presentation}
Present the minimum set of figures and tables needed to support the engineering
conclusion. Each caption should identify the evidence class, operating condition, and
relevant run or test identifier.

\begin{table}[H]
\centering
\caption{Example comparison of analytical and numerical evidence.}
\label{tab:comparison}
\begin{tabularx}{\textwidth}{>{\bfseries}p{0.28\textwidth}XXX}
\toprule
Metric & Analytical & Simulated & Interpretation \\
\midrule
Characteristic frequency & Requirement-derived & Solver-extracted & Agreement must be quantified \\
Primary response & Idealized model & Full numerical model & Differences require physical explanation \\
Efficiency or loss & Simplified estimate & Material and solver dependent & State model fidelity \\
Steering or control & Closed-form prediction & Post-processed result & Report scan loss and asymmetry \\
\bottomrule
\end{tabularx}
\end{table}

\section{Verification checks}
Verification asks whether the equations, scripts, model, and post-processing were
implemented correctly. Useful checks include unit tests, hand calculations,
independent scripts, reciprocity, symmetry, conservation, limiting cases, and
comparison with canonical results.

\section{Validation status}
Validation asks whether the model represents physical reality adequately for the
intended use. Explicitly distinguish model verification from experimental validation.

\begin{measuredevidencebox}
This starter report claims no measured evidence. Replace this box only after a real
prototype, calibrated setup, test procedure, uncertainty statement, and traceable data
record exist.
\end{measuredevidencebox}
